\documentclass[12pt, letterpaper]{article} 
\addtolength{\topmargin}{-30pt}
\usepackage{amsmath, amssymb, amsthm, enumerate, graphicx, fancyhdr}
\pagestyle{fancy}

\begin{document}
\graphicspath{ {./} }

\fancyhead[L]{\textbf{Math 110 HW 1}}
\fancyhead[C]{\textbf{Spring 2026}}
\fancyhead[R]{\textbf{Dani Nguyen}}

\noindent \textbf{Question 1}
\begin{enumerate}[(a)]
    \item $L = (\partial_x + x\partial_y)$
    \begin{align*}
        L(u + v) &= (\partial_x + x\partial_y)(u+v) \\
        &= u_x + v_x + xu_y + xv_y \\
        &= Lu + Lv \\
        L(cu) &= (\partial_x + x\partial_y)(cu) \\
        &= cu_x + cxu_y \\
        &= c(u_x + xu_y) \\
        &= cLu
    \end{align*}
    Therefore $L$ is a linear operator.
    \item $L = (\partial_x + u_y\partial_y)$
    \begin{align*}
        L(u + v) &= (\partial_x + u_y\partial_y)(u+v) \\
        &= u_x + v_x + (u_y)^2 + u_yv_y \\
        &\ne (u_x + (u_y)^2) + (v_x + (v_y)^2) = Lu + Lv 
    \end{align*}
    Therefore, $L$ is NOT a linear operator.
    \item $L = (\partial_x + \partial_y + \frac 1u)$
    \begin{align*}
        L(u + v) &= (\partial_x + \partial_y + \frac 1u)(u+v) \\
        &= u_x + v_x + u_y + v_y + 1 + \frac vu \\
        &\ne (u_x + u_y + 1) + (v_x + v_y + \frac vu) = Lu + Lv
    \end{align*}
    Therefore, $L$ is NOT a linear operator.
    \item $L = (\cos(x)\partial_{xxy} + \arctan(\frac yx)\partial_y)$
    \begin{align*}
        L(u+v) &= (\cos(x)\partial_{xxy} + \arctan(\frac yx)\partial_y)(u+v) \\
        &= \cos(x)u_{xxy} + \arctan(\frac yx)u_y + \cos(x)v_{xxy} + \arctan(\frac yx)v_y \\
        &= Lu + Lv \\
        L(cu) &= (\cos(x)\partial_{xxy} + \arctan(\frac yx)\partial_y)(cu)\\
        &= c\cos(x)u_{xxy} + c\arctan(\frac yx)u_y \\
        &= c(\cos(x)u_{xxy} + \arctan(\frac yx)u_y) \\
        &= cLu
    \end{align*}
    Therefore, $L$ is a linear operator.
\end{enumerate}

\noindent \textbf{Question 2} \\

\noindent \textbf{Question 3}
\begin{enumerate}[(a)]
    \item \begin{align*}
        u_x &= 0 \\
        u &= A(y)
    \end{align*}
    \item \begin{align*}
        u_x &= x^2 + y \\
        u &= \frac 13 x^3
    \end{align*}
\end{enumerate}

\noindent \textbf{Question 4}
\begin{enumerate}[(a)]
    \item 
\end{enumerate}

\noindent \textbf{Question 5}
\begin{enumerate}[(a)]
    \item 
\end{enumerate}

\noindent \textbf{Question 6}

\end{document}